\section{軟弱地面歩行における不安定挙動の特徴}
\label{sec:sim_instability_feature}

本節では，軟弱地面を模擬したシミュレーション結果に基づき，従来手法を用いた歩行において観測された不安定挙動の特徴を整理する．特に，床条件ごとに異なる転倒の発生タイミングおよびその際に生じる力学的特徴に着目し，Case3 および Case4 における転倒要因を分けて考察する．

\subsection{転倒発生タイミングの特徴}
\label{subsec:fall_timing}

シミュレーション結果から，軟弱地面条件においても，床沈み込みそのものが即座に転倒を引き起こすとは限らないことが確認された．Case1 および Case2 では，床沈み込みが発生した場合でも歩行は維持され，転倒には至らなかった．一方で，Case3 および Case4 では転倒が確認されており，それぞれ異なる歩行フェーズにおいて不安定化が生じている．

Case3 では，両脚支持期（Double Support Phase, DS）において急激な床沈み込みが発生し，その直後に姿勢を維持できず転倒に至った．一方，Case4 では，床沈み込み中は姿勢が維持されたものの，両脚支持期から片脚支持期（Single Support Phase, SS）へ遷移する瞬間に転倒が発生した．このことから，軟弱地面歩行における転倒は，単一の要因ではなく，床特性と歩行フェーズの組み合わせによって異なる形で発生することが示唆される．

\subsection{DS--SS 遷移時におけるトルク増大（Case4）}
\label{subsec:torque_increase_case4}

Case4 における転倒挙動を詳細に解析した結果，DS から SS へ切り替わる瞬間において，支持脚足首中心のロール方向の関節トルクが顕著に増大していることが確認された．（グラフ添付予定）

Fig.~\ref{fig:roll_torque_comparison} に示すように，軟弱地面条件では DS--SS 遷移直後に緩やかではあるが大きなトルクが現れており，この過剰トルクの発生が転倒に強く対応している．このことから，Case4 における転倒は，床沈み込みによる外乱が直接的な原因ではなく，歩行フェーズ遷移時に発生する過大な支持脚トルク要求に起因していると考えられる．

\begin{figure}[t]
  \centering
  % \includegraphics[width=0.9\linewidth]{figs/roll_torque_comparison.pdf}
  \caption{剛体床および軟弱地面条件における支持脚ロール方向トルクの比較}
  \label{fig:roll_torque_comparison}
\end{figure}

\subsection{急激な床沈み込みによる姿勢破綻（Case3）}
\label{subsec:instability_case3}

一方，Case3 において観測された転倒は，Case4 とは異なる特徴を示している．Case3 では，床の剛性が低いため，接地直後に急激な床沈み込みが発生し，両脚支持期においてロボットの姿勢が瞬間的に大きく変化した．この急激な姿勢変化により，ZMP が支持脚足裏の支持領域を大きく逸脱し，制御系が姿勢を回復する前に転倒に至った．（グラフ添付予定）

この場合，転倒は DS--SS 遷移に伴うトルク増大ではなく，床沈み込み速度が大きいことによる姿勢条件の急変が主な要因であると考えられる．すなわち，Case3 における不安定化は，従来手法が床沈み込みの時間スケールに対応できないことに起因しており，床特性そのものが歩行制御の前提条件を破綻させている状況であると解釈できる．

\subsection{不安定化要因の整理}
\label{subsec:instability_summary}

以上の結果から，軟弱地面歩行における不安定化には，少なくとも二つの異なる要因が存在することが明らかとなった．一つは，Case3 に見られるような，急激な床沈み込みによる姿勢条件の急変に起因する不安定化である．もう一つは，Case4 に見られるような，床沈み込み後の DS--SS 遷移時における支持脚トルク要求の増大に起因する不安定化である．

本研究では，これらのうち，後者の Case4 における不安定化に着目する．これは，沈み込み中の挙動は安定しているものの，歩行フェーズ遷移時に転倒が発生するという点で，制御手法の改良によって対処可能であると考えられるためである．次章では，この Case4 における不安定化要因を踏まえ，床沈み込みを考慮した歩行制御手法を提案する．
